\chapter{Recover and Preserve Work}\label{ch:recovery}
\section{Recover context, then verify reality}

\begin{headline}
  An interruption can remove the visible conversation while source, remote jobs and
  monitors survive. Recovered context is a lead; current evidence decides.
\end{headline}

Recover the selected harness's native context when supported, then independently
reconcile the chosen brief and state, the owning Git/worktree, the live agent identity
and the relevant remote status. Restore the whole task contract, including downstream
review and delivery obligations \cite[Recover harness context, then verify
  it]{skill:recovery}.

\paragraph{After a reboot, crash or unexplained disconnect:}
\begin{enumerate}
  \item Identify the actual host and time; use a non-interactive preflight for any required
        remote alias.
  \item Inspect host and storage health without cleanup, remounting or service changes.
  \item Locate the selected agent and its native identity using the client or terminal
        multiplexer that owns it.
  \item Read the session records and reconcile them with Git and current job or monitor evidence.
  \item Choose one desired session to resume, record its owner and checkpoint, and continue only
        within its surviving authorization.
\end{enumerate}

\begin{rulebox}[Stop on]
  \begin{itemize}[leftmargin=1.2em]
    \item an unexpected host;
    \item unknown dirty ownership;
    \item disagreement between state and Git;
    \item required credential renewal;
    \item a recovery action that would overwrite information.
  \end{itemize}
  A restart or reconnect does not automatically resume an old autonomous prompt
  \cite[Startup and recovery, steps 1--8; Stop conditions]{skill:recovery}.
\end{rulebox}

\begin{table}[H]
  \small
  \begin{tabularx}{\textwidth}{L{0.26\textwidth}Y}
    \toprule
    Incident & Response \\
    \midrule
    Harness process: a competing or stuck agent process & Identify the exact processes; a
    bounded repair when authorized. This never authorizes replacing the shared SSH
    master. \\
    \rowcolor{ctpSurface0!30}
    SSH transport & The stop-and-inspect rule of Chapter~\ref{ch:remote}. \\
    \bottomrule
  \end{tabularx}
  \caption{Two kinds of incident. Use incident-specific evidence rather than broad
    process-kill patterns \cite[SSH control masters]{skill:remote}.}
  \label{tab:recovery-incidents}
\end{table}

\hook{recovery-clients}

\section{Preserve the result before reducing session scaffolding}

\begin{headline}
  Once operational work settles, promote lasting evidence into an outcome record that lets
  a later reader recover the result without replaying every step.
\end{headline}

The outcome record holds \cite[When]{skill:lifecycle}:
\begin{itemize}
  \item the objective, result and decisions;
  \item report lineage, commits, jobs, reviews and publication;
  \item unresolved limits, with the leftovers it still depends on.
\end{itemize}
Retain unique evidence until a review confirms that consolidation loses nothing essential.

\begin{rulebox}[Supersession]
  Similar report names or dates do not prove supersession: the newer artifact must
  explicitly incorporate or reference the older evidence. A report awaiting human review is
  still active; its source and unique evidence stay protected until the owner's disposition
  is recorded, whether or not the last agent stopped \cite[Review and
    release]{skill:report-production}.
\end{rulebox}

\section{Choose the recovery source for each kind of content}

\begin{headline}
  Git remotes, complementary copies and rebuilds cover different content. Check each
  separately (Table~\ref{tab:recovery}).
\end{headline}

\begin{table}[H]
  \small
  \begin{tabularx}{\textwidth}{L{0.32\textwidth}Y}
    \toprule Content & Recovery evidence \\
    \midrule
    Approved source and records & Reviewed commit actually present on the approved
    writable remote; a local commit alone is not an off-device copy. \\
    \rowcolor{ctpSurface0!30}
    Eligible local-only reading material & Verified complementary copy or documented
    recovery source; respect destination-specific restrictions. \\
    Generated PDF, plots, build state & Preserved inputs, authored source, and a
    reproducible build or generation command. \\
    \rowcolor{ctpSurface0!30}
    Credentials & Excluded from backup; user follows the authority's reissue or
    initialization procedure. \\
    \bottomrule
  \end{tabularx}
  \caption{Recovery classes: content and the evidence that proves it can be restored.}
  \label{tab:recovery}
\end{table}

\begin{choice}{copy tool for complementary copies}
  \begin{itemize}
    \item \recommended\ A tool that uploads additions and replacements without deleting
          remote-only files: the safer routine choice.
    \item \alternative\ A mirror that can delete, only with specific intent.
  \end{itemize}
  Either way, a copy is not Git history and does not provide the same versioned recovery.
\end{choice}

\begin{rulebox}[Restricted material]
  Neither a private Git approval nor a convenient backup route authorizes sending
  restricted material to another service. Verify destination exclusions before copying.
\end{rulebox}

\hook{recovery-backup}

\section{Review cleanup as a separate decision}

\begin{headline}
  An audit finds candidates; it does not authorize removing them. Reversible removals may
  run without per-item approval; everything else waits for the approver.
\end{headline}

A candidate is eligible when it has \cite[Autonomous with
\texttt{-{}-apply}]{skill:lifecycle}:
\begin{itemize}
  \item known ownership and inactivity evidence, and a current inactive status;
  \item no active dependencies;
  \item required publication and a recovery source;
  \item report lineage preserving unique evidence.
\end{itemize}
Age adds a condition, not a substitute: the lifecycle procedure's default minimum is 30
days. A waiver names exact targets and cannot clear a failed semantic gate.

\begin{figure}[H]
  \centering
  \begin{tikzpicture}[node distance=8mm and 14mm]
    \node[work,text width=44mm] (obs) {Expected topology and a dated observation};
    \node[check,text width=44mm,right=of obs] (elig) {Semantic eligibility per candidate};
    \node[work,text width=44mm,below=of obs] (rev) {Reversible: remove without per-item
      approval};
    \node[open,text width=44mm,below=of elig] (appr) {Otherwise: listed for the approver,
      who returns the report hash};
    \node[check,text width=60mm,below=10mm of $(rev.south)!0.5!(appr.south)$] (act)
    {Reclassify each target immediately before acting; refuse drift};
    \draw[flow] (obs)--(elig);
    \draw[flow] (elig.south west)--(rev.north east);
    \draw[flow] (elig)--(appr);
    \draw[flow] (rev.south)--(rev.south |- act.north);
    \draw[flow] (appr.south)--(appr.south |- act.north);
  \end{tikzpicture}
  \caption{Cleanup. This preserves the ability to explain why a target is unused and how its
    important content could be recovered \cite[When; Command]{skill:lifecycle}.}
  \label{fig:recovery-cleanup}
\end{figure}

\paragraph{Holds.} Dirty work, unknown ownership, active references, missing publication,
uncertain lineage, unavailable remote evidence and target drift. Record the specific
missing evidence rather than turning a broad ``looks old'' judgment into approval. A PDF
being ignored by Git does not make it unimportant during review or handoff.

A reversible removal is, for example, a clean worktree whose branch is already on the
remote. Git worktrees are removed through Git, not raw tree deletion \cite[Approval;
  Restore]{skill:lifecycle}.

\begin{rulebox}[No implied authorization]
  Having this handbook, an audit or a dry run authorizes no cleanup action.
\end{rulebox}
